% Part: incompleteness
% Chapter: theories-computability
% Section: intepretability

\documentclass[../../../include/open-logic-section]{subfiles}

\begin{document}

\olfileid[pt]{inc}{tcp}{itp}

\olsection{Teorias nas quais $\Th{Q}$ é interpretável são indecidíveis}

Podemos fortalecer ainda mais esses resultados. Informalmente, uma
interpretação de uma linguagem $\Lang{L_1}$ em outra linguagem
$\Lang{L_2}$ envolve definir o universo, os símbolos de relação e os
símbolos de função de $\Lang{L_1}$ por meio de !!{formula}s em
$\Lang{L_2}$. Embora não nos detenhamos nisso, é possível tornar essa
definição precisa.

\begin{thm}
  Suponha que $\Th{T}$ seja uma teoria em uma linguagem na qual se possa
  interpretar a linguagem da aritmética, de tal maneira que $\Th{T}$ seja
  consistente com a interpretação de $\Th{Q}$. Então $\Th{T}$ é
  indecidível. Se $\Th{T}$ prova a interpretação dos axiomas de
  $\Th{Q}$, então nenhuma extensão consistente de $\Th{T}$ é decidível.
\end{thm}

A prova é apenas uma pequena modificação da prova do último teorema;
poderíamos usar um contraexemplo para obter uma separação de $\Th{Q}$ e
$\Th{\bar Q}$. Podemos tomar $\Th{ZFC}$, a teoria dos conjuntos de
Zermelo--Fraenkel com o axioma da escolha, como um fundamento axiomático
poderoso o suficiente para desenvolver boa parte da matemática usual. Em
$\Th{ZFC}$ podemos definir os números naturais e, por meio dessa
interpretação, os axiomas de $\Th{Q}$ são verdadeiros. Portanto, temos

\begin{cor}
Não há extensão decidível de $\Th{ZFC}$.
\end{cor}

\begin{cor}
Não há extensão completa, consistente e computavelmente !!{axiomatizable} de
$\Th{ZFC}$.
\end{cor}

A linguagem de $\Th{ZFC}$ tem apenas uma relação binária,
$\in$. (Na verdade, nem sequer é preciso ter igualdade.) Portanto, temos

\begin{cor}
A lógica de primeira ordem para qualquer linguagem com um símbolo de relação
binária é indecidível.
\end{cor}

Esse resultado se estende a qualquer linguagem com dois símbolos de função
unária, pois eles podem ser usados para simular um símbolo de relação binária.
Os resultados que acabamos de citar são exatos: verifica-se que a lógica de
primeira ordem para uma linguagem com apenas símbolos de relação
\emph{unária} e no máximo um símbolo de função \emph{unária} é decidível.

Mais uma curiosidade. Sabemos que é indecidível o conjunto das sentenças na
linguagem $\Obj 0$, $'$, $+$, $\times$, $<$ que são verdadeiras no modelo
padrão. Na verdade, podemos definir $<$ em termos dos outros símbolos e,
depois, definir $+$ em termos de $\times$ e $'$. Assim, o conjunto das
sentenças verdadeiras na linguagem $\Obj 0$, $'$, $\times$ é indecidível. Por
outro lado, Presburger mostrou que o conjunto das sentenças na linguagem
$\Obj 0$, $'$, $+$ verdadeiras na linguagem da aritmética é decidível. O
procedimento, porém, é computacionalmente inviável.

\end{document}

